\part{Introduction}

Zeke is a portable operating system for ARM processors.\footnote{Zeke is
portable to at least in a sense of portability between different ARM cores
and architectures.} Figure \ref{figure:zeke} illustrates architectural layers
of the operating system. Scope of Zeke project is to implement a portable
Unix-like\footnote{The project is aiming to implement the core set of kernel
features specified in \ac{POSIX} standard.} operating system from
scratch\footnote{At least almost from scratch, some functionality is derived
from BSD kernels as well as some libraries are taken from other sources.} that
is optimized for ARM architectures and is free of legacy.

In addition to portability another goal has been configurability and adjustable
footprint of the kernel binary as well as the whole OS image. This is achieved
by modular kernel architecture and static compile time configuration. Like for
example \ac{HAL} selection is already almost fully locked at compilation time.

\begin{figure}
  \includegraphics[width=10cm]{pics/zeke}
  \centering
  \caption{Zeke: a Portable Operating System.}
  \label{figure:zeke}
\end{figure}

\section*{List of Abbreviations and Meanings}

\begin{acronym}[UniversalSpacer]
\acro{POSIX}[POSIX]{Portable Operating System Interface for uniX\acroextra{ is a
  family of standards specifying a standard set of features for compatibility
  between operating systems.}}
\acro{vfs}[VFS]{Virtual File System\acroextra{.}}
\acro{dynmem}[dynmem]{(dynamic memory)\acroextra{ is the block memory allocator
  system in Zeke, allocating in 1 MB blocks.}}
\acro{inode}[inode]{\acroextra{inode is the file index node in Unix-style file
  systems that is used to represent a file file system object.}}
\acro{vnode}[vnode]{\acroextra{vnode is like inode but it's used as an
  abstraction level for compatibility between different file systems in Zeke.
  Particularly it's used at VFS level.}}
\acro{ramfs}[ramfs]{RAM file system\acroextra{.}}
\acro{PUnit}[PUnit]{a portable unit testing framework for C\acroextra{.}}
\acro{GCC}[GCC]{GNU Compiler Collection\acroextra{.}}
\acro{GLIBC}[GLIBC]{The GNU C Library\acroextra{.}}
\acro{CPU}[CPU]{Central Processing Unit\acroextra{ is, generally speaking, the
  main processor of some computer or embedded system.}}
\acro{MMU}[MMU]{Memory Management Unit\acroextra{.}}
\acro{DMA}[DMA]{Direct Memory Access\acroextra{ is a feature that allows
 hardware subsystems to commit memory transactions without \acsu{CPU}'s
 intervention.}}
\end{acronym}
